\scanpage{399}\mspage{138}
de $\unsure{\underline{\mathrm{Ty}}|\underline{\mathrm{Tz}}}$ qui soit cofinal dans $\underline{\mathrm{Ty}}$, ce qui signifie
simplement que parmi les $t_n$ il y en a de longueur
utile arbitrairement grande (et comme
\[
n\unsure{\mathrel{\vdash\!\dashv}|\longmapsto}\operatorname{long.ut.}(t_n)
\]
est croissante, cela implique que $\operatorname{long.ut.}(t_n)$
tend vers $+\infty$ pour $n\longrightarrow+\infty$). \deleted{\illegible{1}} Les
faits principaux à vérifier \inserted{ensuite} (si faire se peut)
sont les propriétés \unsure{essentielles} des flèches
dans le diagramme
\diagram{\begin{tikzcd}[ampersand replacement=\&,column sep=small]
\& X \arrow[d,"\tau"] \& \& \\
\& \underline{\mathrm{Ch}}_{\infty}(X)
  \arrow[d,"{\deleted{\unsure{\varphi}}p}"]
  \arrow[ddl,bend right=20,"\sigma"'] \arrow[ddr,bend left=20,"{\beta}",-{Circle[length=1.4pt]}] \& \& \\
\& X\times X \arrow[dl,"\mathrm{pr}_1"'] \arrow[dr,"\mathrm{pr}_2"]
 \& \& \left|\begin{aligned}\sigma\tau&=\beta\deleted{\unsure{\tau}}\tau=\mathrm{id}_X\\
 \unsure{p|\varphi}\tau&=\mathrm{diag}_X\end{aligned}\right.\\
X \& \& X\mathpunct{,} \&
\end{tikzcd}}
et plus généralement, pour une flèche \deleted{\unsure{$X'\longrightarrow X$}}
\[
X'\xrightarrow{f}X
\]
dans $\mathrm{Cat}$, \deleted{\unsure{des}} des flèches dans le diagramme
correspondant
\diagram{\begin{tikzcd}[ampersand replacement=\&,column sep=small]
\& \overline{X'}=\mathrm C(f) \arrow[ddl] \arrow[dr,"g"] \arrow[ddrr,bend left=22,"\bar f"] \& \& \& \bar f\unsure{\alpha|\alpha'}=(\beta g)\alpha=f\\
\& \& \deleted{\unsure{X}}\mathrm{Ch}_{\infty}(X) \arrow[ddl,"\sigma" description] \arrow[dr,"\beta"] \& \& \\
X' \arrow[uur,dashed,bend left=30,"\alpha"] \arrow[dr,"f"'] \& |[draw,ellipse,inner sep=1pt]| \text{cart} \& \& X \& \\
\& X \& \& \&
\end{tikzcd}}
et plus particulièrement de $\beta g$. La flèche
% archive sheet number: 398
